\documentclass[10.5pt,a4paper]{article}
\usepackage[utf8]{inputenc}
\usepackage[T1]{fontenc}
\usepackage[french]{babel}
\usepackage[margin=1.45cm]{geometry}
\usepackage{enumitem}
\usepackage{microtype}
\usepackage[hidelinks]{hyperref}
\setlength{\parindent}{0pt}
\setlength{\parskip}{0pt}
\renewcommand{\baselinestretch}{1.04}
\pagestyle{empty}
\setlist[itemize]{leftmargin=*,nosep,topsep=2pt}
\newcommand{\cvsection}[1]{\par\vspace{5pt}\textbf{\large #1}\par\vspace{1pt}\rule{\linewidth}{0.35pt}\vspace{2pt}}
\begin{document}
\begin{center}
{\Large\textbf{CAMILLE MARTIN}}\\[3pt]
\textbf{Data \& AI Engineer | Python, SQL, Machine Learning | Lyon}\\[3pt]
camille.martin@example.test \;|\; +33 4 00 00 00 01 \;|\; \href{https://example.test/camille}{example.test/camille}
\end{center}

\cvsection{Résumé}
Ingénieure Data \& IA avec deux ans d'expérience sur des produits analytiques, des pipelines cloud et des applications de machine learning. Habituée à relier qualité des données, besoins métier et mise en production de solutions Python/SQL mesurables.

\cvsection{Expérience}
\textbf{GridLab -- Data \& AI Engineer} \hfill 2024--2026
\begin{itemize}
\item Industrialisation de pipelines Python et SQL sur BigQuery pour consolider des données de consommation énergétique, avec contrôles qualité, monitoring et documentation.
\item Développement de modèles de prévision et d'indicateurs opérationnels utilisés par les équipes planification pour analyser demande, anomalies et capacité réseau.
\item Cadrage des besoins, priorisation des livrables et restitution des résultats auprès d'équipes produit, data et opérations dans un environnement Agile.
\end{itemize}
\textbf{UrbanFlow -- Data Analyst, stage} \hfill 2023--2024
\begin{itemize}
\item Construction de datasets SQL réutilisables et de dashboards Looker pour le suivi de mobilité, avec automatisation des contrôles et analyse des écarts.
\item Traduction de questions métier en métriques documentées, analyses pandas et recommandations actionnables.
\end{itemize}

\cvsection{Projets}
\textbf{Prévision de demande énergétique} | Python, pandas, scikit-learn, MLflow
\begin{itemize}
\item Conception d'un pipeline de séries temporelles avec feature engineering, comparaison de modèles, suivi MLflow et analyse d'erreurs par période.
\end{itemize}
\textbf{Assistant RAG pour documentation technique} | Python, FastAPI, embeddings, reranking
\begin{itemize}
\item Développement d'un assistant sur un corpus documentaire avec recherche hybride, réponses sourcées, évaluation et API REST conteneurisée.
\end{itemize}
\textbf{Plateforme temps réel de données mobilité} | Python, Kafka, Spark, Docker
\begin{itemize}
\item Mise en place d'un flux événementiel validant et agrégeant les données, avec tests d'intégration rejouables et métriques de qualité.
\end{itemize}

\cvsection{Compétences}
\textbf{Data Engineering} : Python, SQL, Airflow, dbt, Spark, Kafka, ETL/ELT\\
\textbf{Machine Learning \& GenAI} : scikit-learn, MLflow, RAG, embeddings, évaluation\\
\textbf{Cloud \& Analytics} : GCP, BigQuery, Looker, pandas, qualité des données\\
\textbf{Développement} : FastAPI, APIs REST, Docker, Git, GitHub Actions, CI/CD

\cvsection{Formation}
\textbf{École du Numérique -- Diplôme d'ingénieur, majeure Data \& IA} \hfill 2021--2026\\
Machine learning, statistiques, data engineering, systèmes distribués et cloud.

\cvsection{Langues}
Français natif | Anglais C1 | Espagnol B1

\cvsection{Centres d'intérêt}
Énergie | Mobilité | Photographie
\end{document}
